% !Mode:: "TeX:UTF-8"
\chapter{总结与展望}
\label{cha:conclusion}

\section{本文工作总结}
本文围绕 ``复杂环境下无人船的博弈策略研究'' 这一核心课题, 以欠驱动水面无人艇为研
究对象, 结合生物启发、威胁势场理论与深度强化学习方法, 系统研究了单艇追逃博弈、
多艇协同围捕博弈和 TAD 目标防御博弈三类典型博弈任务的策略优化问题. 本文的主要
研究工作可以总结为以下三个方面:

(1) \textbf{针对三维复杂环境下欠驱动无人艇追逃博弈问题}, 基于近端策略优化算法提
出了一种融合三段式观测空间 (运动观测、环境感知观测与扰动观测) 与鲨鱼捕食行为启
发奖励机制的追逃策略. 扰动观测的显式建模增强了算法对环境不确定性的适应能力; 生
物启发式奖励函数缓解了传统奖励稀疏导致的训练困难. 仿真实验表明, 所提算法在不同
扰动强度与典型追逃场景下均表现出良好的追击性能与鲁棒性.

(2) \textbf{针对复杂海洋环境下多无人艇协同围捕博弈问题}, 首先, 基于六自由度非线性
运动学模型, 给出了协同围捕成功的严格数学判据, 填补了现有文献在围捕评价标准上的空
白; 其次, 提出能量波峰 (EC) 与能量波谷 (ET) 概念, 将威胁势场转化为反映接近-成型-保
持三阶段的密集奖励; 再次, 在 MAPPO 的 Critic 网络中引入多头注意力模块, 构建了
MHD-MAPPO 算法, 显著提升了多艇协同场景下的收敛效率与决策质量; 最后, \textbf{搭
建了多无人艇物理实验平台, 完成了真实水面环境下的 2 对 1 与 3 对 1 协同围捕实物实
验}, 系统验证了所提算法从仿真到实物的迁移能力和工程可行性.

(3) \textbf{针对非自杀式防御者的 TAD 目标防御博弈问题}, 建立了动态目标护航与静态
目标固守的两阶段任务模型; 在策略网络中引入残差自注意力与循环神经结构以实现智能
体间关联建模与时序感知; 通过扩展卡尔曼滤波实现进攻者状态预测, 使防御者具备前瞻
性决策能力; 受鲸群防御行为启发设计势场化奖励函数. 仿真与消融实验表明, 所提 TP-
POCA 算法在动态和静态目标场景下均显著优于主流多智能体强化学习基线方法.

\section{研究展望}
尽管本文在复杂环境下无人艇博弈策略研究方面取得了一定成果, 但仍存在值得进一步深
入研究的问题, 主要包括以下几个方面:

(1) \textbf{异构多无人艇博弈策略研究.} 本文研究的多无人艇博弈主要面向同构无人艇,
在实际海洋任务中, 不同吨位、不同能力的异构无人艇往往需要协同执行任务. 未来可研
究异构无人艇间能力互补与角色分工的博弈策略.

(2) \textbf{面向真实海况的大范围物理实验.} 虽然本文在第四章完成了协同围捕的水池
及湖面实物实验, 但尚未在大风浪、强流速等真实海况下验证算法的工程可靠性. 未来可
扩展至近海海域完成大范围、多艇数的实海实验.

(3) \textbf{空-海跨域协同博弈研究.} 随着无人机、无人潜航器等平台的快速发展, 空-海-
潜跨域协同博弈正成为新一代海上作战关键技术. 未来研究可拓展到 USV 与 UAV / UUV
之间的协同博弈, 实现立体化的博弈对抗.

(4) \textbf{基于大模型的无人艇自主决策研究.} 近年来大语言模型与多模态基础模型在决
策与规划中展现出巨大潜力. 将大模型与强化学习相结合, 实现面向未知任务的零样本或
少样本博弈策略学习, 是下一步值得探索的方向.

(5) \textbf{博弈策略的安全性与可解释性研究.} 在实际应用中, 无人艇博弈策略的安全边
界与决策可解释性至关重要. 未来可结合可信人工智能、形式化验证等方法, 进一步提升
策略的安全性与透明度.
